% Szakdolgozat sablon — fejezet- és alfejezet-struktúra a
% admin/szakdolgozat/szakdolgozat_felepites_vazlat.pdf tartalma szerint (EKF / Kovács Ádám vázlat).
% Osztály: thesis-ekf (lásd thesis-ekf-template-en.tex); a thesis-ekf.cls elérhetőnek kell lennie.
% Fordítás: pdflatex szakdolgozat-sablon.tex (kétszer, ha szükséges)
\documentclass{thesis-ekf}
\usepackage[T1]{fontenc}
\usepackage[utf8]{inputenc}
\usepackage[magyar]{babel}
\usepackage{booktabs}
\usepackage{graphicx}
\usepackage{hyperref}
\hypersetup{colorlinks=true,linkcolor=blue,urlcolor=blue}

\begin{document}

\institute{Matematika- és Informatikai Intézet}
\title{Szakdolgozat címe}
\author{Hallgató neve\\szak}
\supervisor{Konzulens neve\\beosztás}
\city{Eger}
\date{2026}
\maketitle

\tableofcontents
\listoffigures
\listoftables

\chapter*{Formai követelmények és szerkesztési útmutató}
\addcontentsline{toc}{chapter}{Formai követelmények és szerkesztési útmutató}

A szakdolgozatot \LaTeX{} rendszerben kell megírni. A leggyakrabban használt sablon Tomács Tibor oktatóanyagai és mintái alapján készül; elérhetőségek és telepítőcsomag a tanszéki oldalon (\url{https://tomacstibor.uni-eszterhazy.hu/szakdolgozat.html}).

A dolgozat szövegének stílusa törekedjen az egyensúlyra: nem baráti hangvételű, de nem is teljesen tudományosan távolságtartó. A cél egy közérthető, szakmailag igényes, de olvasmányos szöveg.

\begin{table}[htbp]
\centering
\caption{A legfontosabb formai követelmények összefoglalása}
\begin{tabular}{@{}ll@{}}
\toprule
\textbf{Formai elem} & \textbf{Előírás} \\
\midrule
Tartalomjegyzék & \texttt{\string\tableofcontents} \\
Ábra- és táblázatjegyzékek & \texttt{\string\listoffigures}, \texttt{\string\listoftables} \\
Ábrafeliratok & Számozott, forrással és megtekintési dátummal \\
Irodalomjegyzék & Számozott, IEEE-stílusban ajánlott \\
Nyilatkozat & A dolgozat végén PDF formátumban csatolva \\
\bottomrule
\end{tabular}
\end{table}

Leadási szabályok: a végleges anyag feltöltése a Neptunon keresztül történik; a PDF ne legyen fájlvédett; a fájlnév a Neptun-kód legyen (nagybetűvel); ajánlott a forrás és kapcsolódó állományok \texttt{.zip} archívuma. A részletes feltételeket a mindenkori Tanulmányi és vizsgaszabályzat és a tanszéki tájékoztató tartalmazza.

\chapter{Bevezetés}

\section{Témaválasztás megindoklása}
\begin{itemize}
  \item Miért pont ezt a témát választottad?
  \item Milyen aktuális iparági trendek, szakmai kihívások vagy személyes célok kapcsolódnak hozzá?
\end{itemize}

\section{Téma bemutatása}
\begin{itemize}
  \item Rövid áttekintés a témáról, releváns szakirodalmi háttér és alapfogalmak.
  \item Kapcsolódó tudományos cikkek, piaci elemzések, ipari megoldások.
\end{itemize}

\section{Cél leírása}
\begin{itemize}
  \item Mit szeretnél elérni a dolgozattal?
  \item Milyen várható eredményekre, fejlesztésekre számítasz?
  \item Főbb mérföldkövek, amelyek a megvalósítás lépéseit jelölik.
\end{itemize}

\section{Hipotézis}
Nem szükséges teljesen új dolgot felfedezni; egy ismert módszer vagy algoritmus adaptálása is megfelelő lehet.
\begin{itemize}
  \item Példa: „Az általam alkalmazott algoritmus gyorsabb lesz, mint a korábbi változatok.''
  \item Példa: „Az új fejlesztési módszertannal csökkenthető a hibák száma.''
\end{itemize}

\begin{sloppypar}
\textit{Megjegyzés:} A hipotézis megfogalmazása elsősorban akkor szükséges, ha a dolgozat iránya kutatásalapú, vagy ha mérhető teljesítményjavulást, új fejlesztési irányt, optimalizációt kívánsz igazolni. Ha a téma ezt nem indokolja, a célkitűzés rész megfelelően helyettesíti.
\end{sloppypar}

\section{Hasonló alkalmazások}
\begin{itemize}
  \item Milyen hasonló megoldások léteznek már a szakirodalomban vagy a piacon?
  \item Hiányosságok, lehetséges továbbfejlesztések.
  \item Miben tud segíteni az általad fejlesztett alkalmazás, miben új vagy jobb?
\end{itemize}

\section{Hivatkozások és bemutató linkek}
A bevezetés végén javasolt három hivatkozás elhelyezése, ha a projekt szoftveres megvalósítása online is bemutatható. Az első link a forráskódra, a második a buildre (webes alkalmazás esetén weboldal), a harmadik egy 3--5 perces bemutató videóra mutasson.
\begin{itemize}
  \item GitHub (forráskód): \url{https://github.com/felhasznalo/projektKod}
  \item GitHub (build): \url{https://github.com/felhasznalo/projektBuild}
  \item Videó (rövid projektbemutató): \url{https://www.youtube.com/}
\end{itemize}

\chapter{Technológiai áttekintés}

\section{Fejlesztői környezetek}
\begin{itemize}
  \item Milyen IDE-t használsz (pl.\ PyCharm, Visual Studio Code, IntelliJ stb.)?
  \item Fontosabb funkciók, bővítmények, plug-inek.
  \item Esetleges hátrányok, nehézségek.
\end{itemize}

\section{Programozási nyelvek}
\begin{itemize}
  \item Miért pont ezt (ezeket) a nyelv(ek)et választottad?
  \item Előnyök, korlátok, közösségi támogatás.
\end{itemize}

\section{Verziókezelés}
\begin{itemize}
  \item Verziókövető eszköz (pl.\ Git) használata.
  \item Branch-struktúra, Git-flow, egyéb módszerek.
  \item CI/CD megoldások (automatizált teszt, build).
\end{itemize}

\section{Használt csomagok és keretrendszerek}
Ebben az alfejezetben mutasd be a dolgozatban alkalmazott fontosabb könyvtárakat, csomagokat vagy fejlesztési keretrendszereket:
\begin{itemize}
  \item Milyen problémák megoldására használod őket?
  \item Melyik nyelvhez kapcsolódnak?
  \item Milyen alternatívák vannak, és miért pont ezt választottad?
  \item Van-e közösségi támogatás, dokumentáció, hosszú távú karbantartás?
\end{itemize}

\chapter{Rendszerterv}

\section{A rendszer célja}
\begin{itemize}
  \item Pontosan milyen problémát old meg a rendszer?
  \item Célfelhasználók, felhasználási forgatókönyvek.
\end{itemize}

\section{Követelmények listája}
\textbf{Funkcionális követelmények}
\begin{itemize}
  \item Adatbevitel, riportok, statisztika stb.
  \item Szerepkörök (pl.\ felhasználó, admin).
\end{itemize}
\textbf{Nem-funkcionális követelmények}
\begin{itemize}
  \item Teljesítmény, skálázhatóság.
  \item Biztonság, autentikáció, titkosítás.
  \item Felhasználói élmény (UI, reszponzivitás).
\end{itemize}

\section{Használati esetek (Use case-ek)}
\begin{itemize}
  \item Use case diagramok.
  \item Forgatókönyvek (például „Felhasználó regisztrál'', „Admin riportot készít'').
\end{itemize}

\section{Adatbázis-terv}
Amennyiben a dolgozat tartalmaz adatbázist, elvárás, hogy az legalább a harmadik normálformának (3NF) megfeleljen. Be kell mutatni a táblák közötti kapcsolatokat, az attribútumfüggőségeket, kulcsokat, relációkat, valamint ajánlott ER-diagramot is mellékelni.

\section{UML-diagramok, folyamatábrák, struktogramok}
\begin{itemize}
  \item UML osztálydiagram (Class diagram).
  \item UML szekvenciadiagram (pl.\ bejelentkezési folyamat).
  \item Folyamatábrák, struktogramok, ha szükséges.
\end{itemize}

\chapter{Saját szoftver megvalósítása}

\section{Az alkalmazás felépítése}
\begin{itemize}
  \item Átfogó architektúra (frontend, backend, adatbázis).
  \item Rétegek és komponensek, miként kommunikálnak egymással.
\end{itemize}

% Példa ábra — a \caption generálja a jegyzékbejegyzést
\begin{figure}[htbp]
\centering
\fbox{\parbox{0.85\textwidth}{\centering\vspace{3em}\textit{[Ábra: architektúra vagy képernyőkép helye]}\vspace{3em}}}
\caption{Példa ábra felirattal és forrásmegjelöléssel}
\end{figure}

\section{Backend bemutatása}
\begin{itemize}
  \item Keretrendszer (Django, Flask, Spring Boot stb.).
  \item REST API, adatkezelés, külső szolgáltatások, háttérfolyamatok.
\end{itemize}

\section{Frontend bemutatása}
\begin{itemize}
  \item Framework (Angular, React, Vue stb.).
  \item Komponensek, oldalak, navigáció.
  \item UI/UX megfontolások (reszponzív design, validáció).
\end{itemize}

\section{Modulok bemutatása}
A szoftveres megvalósítást érdemes modulok halmazaként strukturálni; minden modulnál:
\begin{itemize}
  \item \textbf{Cél:} feladatok, illeszkedés az alkalmazás többi részéhez.
  \item \textbf{Bemenetek, kimenetek, funkciók.}
  \item \textbf{Rövid kódrészlet} (legfeljebb háromnegyed oldal), GitHub-hivatkozás hosszabb részekhez.
  \item \textbf{UML/folyamatábra} (opcionális) összetett logikához.
  \item \textbf{Függőségek:} kapcsolat más modulokkal (közös adatbázis, API-hívások stb.).
\end{itemize}

\section{Futtatási leírás}
\begin{itemize}
  \item Rendszerkövetelmények (operációs rendszer, függőségek, könyvtárak).
  \item Telepítési lépések (\texttt{git clone}, \texttt{npm install}, \texttt{pip install -r requirements.txt} stb.).
  \item Adatbázis inicializálás, konfigurációs beállítások, migrációk.
  \item Indítás (pl.\ \texttt{npm start}, \texttt{php artisan serve}).
\end{itemize}

\chapter{Tesztelés}

\section{Manuális tesztelés}
\begin{itemize}
  \item Tesztforgatókönyvek, rövid leírás, bemenet, kimenet, eredmények.
  \item Hibajavítás, újra tesztelés dokumentálása.
\end{itemize}

\section{Automatikus tesztelés}
\begin{itemize}
  \item Használt keretrendszerek (JUnit, pytest, PHPUnit stb.).
  \item Tesztlefedettség (coverage-jelentések).
  \item CI integráció (GitHub Actions, Jenkins), pipeline menete.
\end{itemize}

\chapter{Konklúzió és eredmények összehasonlítása}
\begin{itemize}
  \item A bevezetőben (Hipotézis) leírtak igazolása vagy cáfolata.
  \item Összevetés a szakirodalomban ismert eredményekkel.
  \item Esetleges mérési adatok kiértékelése, grafikonok, táblázatok.
\end{itemize}

\chapter{Összegzés}

\section{Eredmények összefoglalása}
\begin{itemize}
  \item Mely célkitűzések valósultak meg?
  \item Tapasztalatok a fejlesztés során.
\end{itemize}

\section{Tanulságok}
\begin{itemize}
  \item Milyen problémákkal szembesültél, hogyan oldottad meg őket?
  \item Milyen új tudásra tettél szert?
\end{itemize}

\section{Jövőbeli fejlesztések}
\begin{itemize}
  \item Lehetséges továbbfejlesztések, funkcióbővítések.
  \item Skálázhatóság, integrációk, optimalizálási lehetőségek.
\end{itemize}

\textit{Megjegyzés:} A dolgozat ne tartalmazzon hosszú forráskód-részleteket; javasolt az architektúra ábrákkal, folyamatábrákkal, UML-diagramokkal történő szemléltetése. A teljes kód legyen elérhető GitHub linken.

\begin{thebibliography}{9}
\addcontentsline{toc}{chapter}{\bibname}
\bibitem{ekf-szakd}
Eszterházy Károly Katolikus Egyetem, „Szakdolgozati sablon, bírálatilap és további információk,'' \url{https://uni-eszterhazy.hu/matinf/szakdolgozat}.
\bibitem{w3c-html}
World Wide Web Consortium, „HTML5 Specification,'' \url{https://www.w3.org/TR/html5/}.
\bibitem{mdn-css}
Mozilla Developer Network, „CSS: Cascading Style Sheets,'' \url{https://developer.mozilla.org/en-US/docs/Web/CSS}.
\end{thebibliography}

\chapter*{Nyilatkozat csatolása}
\addcontentsline{toc}{chapter}{Nyilatkozat csatolása}

A végleges PDF-hez a tanszék által közzétett nyilatkozatot kell csatolni. Általános lépések:
\begin{enumerate}
  \item A \texttt{nyilatkozat/nyilatkozat.tex} sablon kitöltése és PDF-be fordítása.
  \item Nyomtatás, aláírás, beszkenelés PDF-be.
  \item A szakdolgozat végén a \texttt{pdfpages} csomaggal: \texttt{\string\includepdf\{nyilatkozat/nyilatkozat.pdf\}} az \texttt{\string\end\{document\}} előtt.
\end{enumerate}

% \usepackage{pdfpages}
% \includepdf{nyilatkozat/nyilatkozat.pdf}

\end{document}
